\begin{ex}
Một lớp có 32 học sinh, trong đó $\frac{1}{4}$ số học sinh là học sinh giỏi. Hỏi lớp đó có bao nhiêu học sinh giỏi?
\shortans{8}
\loigiai{
$32 : 4 = 8$ học sinh.
}
\end{ex}

\begin{ex}
Một khúc gỗ dài 24 m. Người ta cưa lấy $\frac{1}{3}$ khúc gỗ. Hỏi đoạn gỗ được cưa lấy dài bao nhiêu mét?
\shortans{8}
\loigiai{
$24 : 3 = 8$ m.
}
\end{ex}

\begin{ex}
Một giỏ có 35 quả cam. Mẹ lấy ra $\frac{1}{5}$ số cam để vắt nước. Hỏi mẹ đã lấy ra bao nhiêu quả cam?
\shortans{7}
\loigiai{
$35 : 5 = 7$ quả cam.
}
\end{ex}

\begin{ex}
Một ngày có 24 giờ. Một phần tư ($\frac{1}{4}$) của một ngày là bao nhiêu giờ?
\shortans{6}
\loigiai{
$24 : 4 = 6$ giờ.
}
\end{ex}

\begin{ex}
Một giờ có 60 phút. Một phần sáu ($\frac{1}{6}$) của một giờ là bao nhiêu phút?
\shortans{10}
\loigiai{
$60 : 6 = 10$ phút.
}
\end{ex}

\begin{ex}
Một sợi dây dài 42 cm. Hỏi $\frac{1}{7}$ sợi dây đó dài bao nhiêu xăng-ti-mét?
\shortans{6}
\loigiai{
$42 : 7 = 6$ cm.
}
\end{ex}

\begin{ex}
Một hộp có 48 viên bi. Nam lấy $\frac{1}{8}$ số bi trong hộp. Hỏi Nam lấy được bao nhiêu viên bi?
\shortans{6}
\loigiai{
$48 : 8 = 6$ viên bi.
}
\end{ex}

\begin{ex}
Đàn gà có 54 con, trong đó $\frac{1}{9}$ là gà trống. Hỏi có bao nhiêu con gà trống?
\shortans{6}
\loigiai{
$54 : 9 = 6$ con gà trống.
}
\end{ex}

\begin{ex}
Một cuộn vải dài 36 m. Người ta đã bán đi $\frac{1}{6}$ cuộn vải. Hỏi người ta đã bán bao nhiêu mét vải?
\shortans{6}
\loigiai{
$36 : 6 = 6$ m.
}
\end{ex}

\begin{ex}
Tìm $\frac{1}{2}$ của 18 kg.
\shortans{9}
\loigiai{
$18 : 2 = 9$ kg.
}
\end{ex}

\begin{ex}
Tìm $\frac{1}{3}$ của 27 lít.
\shortans{9}
\loigiai{
$27 : 3 = 9$ lít.
}
\end{ex}

\begin{ex}
Tìm $\frac{1}{7}$ của 49 chiếc bút.
\shortans{7}
\loigiai{
$49 : 7 = 7$ chiếc.
}
\end{ex}

\begin{ex}
Tìm $\frac{1}{8}$ của 64 bạn học sinh.
\shortans{8}
\loigiai{
$64 : 8 = 8$ bạn.
}
\end{ex}

\begin{ex}
Tìm $\frac{1}{9}$ của 72 quả bưởi.
\shortans{8}
\loigiai{
$72 : 9 = 8$ quả.
}
\end{ex}

\begin{ex}
Một cửa hàng có 40 bao gạo, đã bán $\frac{1}{5}$ số bao đó. Hỏi cửa hàng còn lại bao nhiêu bao gạo?
\shortans{32}
\loigiai{
Đã bán: $40 : 5 = 8$ bao. Còn lại: $40 - 8 = 32$ bao gạo.
}
\end{ex}

\begin{ex}
Một rổ táo có 24 quả. Hà ăn $\frac{1}{4}$ số táo đó. Hỏi rổ táo còn lại bao nhiêu quả?
\shortans{18}
\loigiai{
Hà ăn: $24 : 4 = 6$ quả. Còn lại: $24 - 6 = 18$ quả.
}
\end{ex}

\begin{ex}
Bác nông dân thu được 63 kg bí ngô và đã bán đi $\frac{1}{7}$ số bí ngô đó. Hỏi bác còn lại bao nhiêu ki-lô-gam bí ngô?
\shortans{54}
\loigiai{
Đã bán: $63 : 7 = 9$ kg. Còn lại: $63 - 9 = 54$ kg.
}
\end{ex}

\begin{ex}
Có 81 bông hoa xếp vào các bình, mỗi bình có $\frac{1}{9}$ số hoa. Hỏi mỗi bình có bao nhiêu bông hoa?
\shortans{9}
\loigiai{
$81 : 9 = 9$ bông hoa.
}
\end{ex}

\begin{ex}
Một mảnh vườn hình vuông có chu vi 28 m. Độ dài mỗi cạnh của mảnh vườn bằng $\frac{1}{4}$ chu vi. Hỏi cạnh mảnh vườn dài bao nhiêu mét?
\shortans{7}
\loigiai{
$28 : 4 = 7$ m.
}
\end{ex}

\begin{ex}
Tìm một số biết rằng $\frac{1}{6}$ của số đó bằng 7.
\shortans{42}
\loigiai{
Số đó là $7 \times 6 = 42$.
}
\end{ex}

\begin{ex}
Tìm một số biết rằng $\frac{1}{8}$ của số đó bằng 5.
\shortans{40}
\loigiai{
Số đó là $5 \times 8 = 40$.
}
\end{ex}

\begin{ex}
Tìm một số biết rằng $\frac{1}{9}$ của số đó bằng 9.
\shortans{81}
\loigiai{
Số đó là $9 \times 9 = 81$.
}
\end{ex}

\begin{ex}
Một bao gạo nặng 50 kg. Sau khi dùng $\frac{1}{5}$ bao gạo thì còn lại bao nhiêu ki-lô-gam?
\shortans{40}
\loigiai{
Đã dùng: $50 : 5 = 10$ kg. Còn lại: $50 - 10 = 40$ kg.
}
\end{ex}

\begin{ex}
Đội văn nghệ có 20 bạn, trong đó $\frac{1}{4}$ là bạn nam. Hỏi đội văn nghệ có bao nhiêu bạn nữ?
\shortans{15}
\loigiai{
Số bạn nam: $20 : 4 = 5$ bạn. Số bạn nữ: $20 - 5 = 15$ bạn.
}
\end{ex}

\begin{ex}
Biết $\frac{1}{3}$ của số $A$ là 8. Hỏi số $A$ là bao nhiêu?
\shortans{24}
\loigiai{
Số $A$ là $8 \times 3 = 24$.
}
\end{ex}
